% Budget: ca. 0,4 Seiten. Verantwortlich: Person A
\section{Hintergrund}
\label{sec:hintergrund}

% Definition KI-Coding-Agent: LLM-basiert, mehrschrittig, nutzt Werkzeuge (Dateien, Shell, Tests).
% Abgrenzung zu reiner Code-Completion.
\todo{Definition KI-Coding-Agent und Abgrenzung}

% Autonomiestufen: von Vorschlägen über Human-in-the-Loop bis vollautonom.
% Begriff „Workflow“ für dieses Paper festlegen.
\todo{Autonomiestufen und Workflow-Begriff}

% Kurzer Überblick über Agentensysteme und Benchmarks, z.B. SWE-bench~\cite{jimenez2024swebench},
% SWE-agent~\cite{yang2024sweagent}, OpenHands~\cite{wang2025openhands} und Agentless~\cite{xia2024agentless}
\todo{Überblick über Agentensysteme und Benchmarks~\cite{jimenez2024swebench,yang2024sweagent,wang2025openhands,xia2024agentless}}
